\documentclass[10pt,a4paper]{amsart}
\usepackage[margin=19mm]{geometry}
\usepackage{amsmath,amssymb,mathtools}
\usepackage[scr=rsfs,cal=euler]{mathalfa}
\usepackage{xfrac}
\usepackage[unicode,hidelinks,bookmarks=false]{hyperref}
\newcommand{\ie}{i.e.\ }
\newcommand{\cf}{cf.\ }
\newcommand{\resp}{respectively\ }
\newcommand{\Subsection}{\subsection}
\providecommand{\nobreakdash}{\nobreak\hbox{-}}
\setlength{\emergencystretch}{2em}
\multlinegap0pt
\usepackage{bm}
\usepackage{bbm}
\usepackage{dsfont}
\usepackage{xspace}
\usepackage{cancel}
\usepackage{mathtools}
\usepackage{amsthm}
\makeatletter
\@ifundefined{theorem}{\newtheorem{theorem}{Theorem}}{}
\@ifundefined{lemma}{\newtheorem{lemma}[theorem]{Lemma}}{}
\makeatother
\title[The moments of split greatest common divisors]{The moments of split greatest common divisors}
\author{Abhishek Jha, Ayan Nath, Emanuele Tron}
\date{Source: arXiv:2408.05820v3 (CC BY 4.0)}
\begin{document}
\maketitle

\noindent\emph{Source and licence.} The moments of split greatest common divisors. Version arXiv:2408.05820v3 (CC BY 4.0), distributed under the Creative Commons Attribution 4.0 International licence, as recorded in the arXiv licence metadata for this exact version and shown on the abstract page https://arxiv.org/abs/2408.05820v3. Full rights notice: licenses/arxiv-2408.05820-CC-BY-4.0.txt.\par\smallskip

\noindent\textit{Editorial scope.} Counted unit: the Lemma whose source label is \texttt{lemma 2.4} in the article (source0.tex lines 537--539, source label \texttt{lemma 2.4}), delivered together with the article's own complete original proof of it (source0.tex lines 540--557, one \texttt{proof} environment of 18 lines). No mathematics is written, repaired or re-derived: statement and proof are the article's own text line for line. Required context retained uncounted: none. Every definition, display and label of those lines is retained unchanged. Omitted from this package and not redistributed: 1--536, 558--654. Nothing inside a retained excerpt is omitted or altered: every retained body is delivered exactly as the article's lines are written, so no omission receipt of any kind accompanies this package. Citations: the retained excerpts carry no citation command at all, so no bibliography entry of the article is transcribed and no reference is redistributed. Reference adaptation: none was needed and none was made; every reference inside the delivered unit resolves inside it, so no unresolved reference or citation is produced. Printed slips kept literal: no printed word, symbol, hypothesis or number of the counted statement or of its proof was altered. Layout and class: the article's own class is \texttt{amsart} with packages including amsmath, colonequals, times, fullpage, amssymb, amsthm, geometry, microtype, color, graphicx, hyperref, enumitem, dsfont; the wrapper is the lane's pinned \texttt{amsart} editorial wrapper at 10pt on A4, which restates the theorem-like environments the retained text needs and adds the article's own macro definitions that the retained text needs (none), with extra wrapper packages (bm, bbm, dsfont, xspace, cancel, mathtools). The delivered numbering is the unit's own. Assets: no retained span contains a figure environment, a graphics-inclusion command, a PDF-page inclusion, a table or a TikZ picture; no image file is copied into this package, the package declares no asset dependency and no image file is redistributed with it. Licence: version arXiv:2408.05820v3 of this article is distributed under the Creative Commons Attribution 4.0 International licence, as recorded in the arXiv licence metadata for this exact version and shown on the abstract page https://arxiv.org/abs/2408.05820v3. A full rights notice is delivered at licenses/arxiv-2408.05820-CC-BY-4.0.txt. Characters: every retained excerpt is pure ASCII, so the delivered engine, XeLaTeX with its default Latin Modern text font, reports no missing character.
\begin{lemma}\label{lemma 2.4}
    Let $d$ be a squarefree integer. Then the number of squarefree $n\le x$ such that $\gcd(n,a_2)=1$ and $d\mid z(n)$ is at most \[\omega(d)^{\omega(d)} \,\frac{x}{d} (C (\log_{2}x)^2)^{\omega(d)}\] for some constant $C>0$.
\end{lemma}

\begin{proof}
    Let \[\varphi^{*}(n)=\prod_{p\mid n} \left(p - (-1)^{p-1}\eta(p)\right)\] be a modified Euler function, where $\eta(p)=(\Delta/p)$ when $p\neq 2$ and $\eta(2)= \mathds{1}_{2\nmid \Delta}$. Clearly, $z(n) \mid \varphi^{*}(n)$ for all squarefree $n$. 
    
Since $d$ divides $\varphi^{*}(n)$, we can write $d = \prod_{p \mid n} a_p$, where each $a_p$ divides $p - (-1)^{p-1}\eta(p)$. Discarding those $a_p = 1$, we see that $n$ induces a (not necessarily unique) factorization of $d$, whereby by a \emph{factorization} of $d$ we mean a decomposition of $d$ as a product of factors strictly larger than $1$, where the order of the factors is not taken into account. For each possible factorization of $d$, we estimate the number of $n \leq x$ as in the lemma statement which induce this factorization.

Let $d = a_1 a_2 \cdots a_k$ be a factorization of $d$. If $n$ induces this factorization, then there are distinct primes $p_1, \ldots, p_k$ dividing $n$ with $p_i \equiv -1\text{ or }0\text{ or }1 \pmod{a_i}$ for each $1 \leq i \leq k$. So, the number of such $n \leq x$ is bounded above by 
\[x\prod_{i=1}^{k}\left(\sum_{\substack{p_i\le x\\ p_i\equiv -1,0,1\bmod{a_i}}}\frac1{p_i}\right)\]
For each $i$, the inner sum is bounded by $\ll \frac{\log_{2} x}{\varphi(a_i)}+\frac1{a_i}\ll \frac{(\log_{2} x)^{2}}{a_i}$. Thus, for an absolute constant $C>0$, the number of $n$ is at most 
$$
x \, \prod_{i=1}^k \frac{C (\log_2 x)^{2}}{a_i} 
= \frac{x}{d} (C (\log_2 x)^{2})^k 
\leq \frac{x}{d} (C (\log_2 x)^{2})^{\omega(d)}.$$ (The last inequality uses the observation that each factorization of $d$ involves at most $\omega(d)$
factors.) Since $d$ is squarefree, the number of factorizations of $d$ is given by $B_{\omega(d)}$
, where $B_l$
(the $l^{\text{th}}$ Bell number ) stands for the number of set-partitions of an $l$-element set.
Since any partition of an $l$-element set involves at most $l$ components, we always
have $B_l\le  l^{l}$. Taking $l = \omega(d)$ completes the proof.
\end{proof}

\end{document}
